\documentclass[10pt,landscape]{article}
\usepackage[utf8]{inputenc}
\usepackage[spanish]{babel}
\usepackage{multicol}
\usepackage{geometry}
\usepackage{amsmath, amssymb}
\usepackage{booktabs}
\usepackage{xcolor}
\usepackage{tcolorbox}

% Configuración de página
\geometry{
    paper=a4paper,
    left=0.5cm,
    right=0.5cm,
    top=0.5cm,
    bottom=0.5cm,
    landscape
}

% Colores
\definecolor{azulprofundo}{RGB}{0, 56, 147}
\definecolor{azulclaro}{RGB}{41, 128, 185}
\definecolor{verdeexito}{RGB}{39, 174, 96}
\definecolor{grisclaro}{RGB}{236, 240, 241}

% Configuración de secciones
\setlength{\columnsep}{0.5cm}
\setlength{\columnseprule}{0.2pt}

% Comandos personalizados
\newcommand{\comando}[1]{\texttt{\textcolor{azulprofundo}{#1}}}
\newcommand{\codigo}[1]{\texttt{\textcolor{azulclaro}{#1}}}

\begin{document}

\begin{center}
    {\Huge \textbf{Cheatsheet \LaTeX}}
    \vspace{0.2cm}
    
    {\large Referencia rápida de comandos para estudiantes de ingeniería}
    \vspace{0.5cm}
    
    \rule{0.8\textwidth}{0.5pt}
\end{center}

\begin{multicols}{3}

% ============================================
% ESTRUCTURA DEL DOCUMENTO
% ============================================
\tcbox[colback=grisclaro, colframe=azulprofundo, arc=3pt, boxsep=2pt, left=2pt, right=2pt, top=2pt, bottom=2pt]{
    \textbf{\textcolor{azulprofundo}{📄 ESTRUCTURA DEL DOCUMENTO}}
}
\vspace{0.3cm}

\begin{tabular}{@{}ll@{}}
    \toprule
    \textbf{Comando} & \textbf{Función} \\
    \midrule
    \comando{\textbackslash documentclass\{...\}} & Tipo de documento \\
    \comando{\textbackslash usepackage\{...\}} & Cargar paquete \\
    \comando{\textbackslash begin\{document\}} & Inicio del cuerpo \\
    \comando{\textbackslash end\{document\}} & Fin del cuerpo \\
    \comando{\textbackslash title\{...\}} & Título del documento \\
    \comando{\textbackslash author\{...\}} & Autor \\
    \comando{\textbackslash date\{...\}} & Fecha \\
    \comando{\textbackslash maketitle} & Genera la portada \\
    \bottomrule
\end{tabular}

\vspace{0.5cm}

% ============================================
% CLASES DE DOCUMENTO
% ============================================
\tcbox[colback=grisclaro, colframe=azulprofundo, arc=3pt, boxsep=2pt, left=2pt, right=2pt, top=2pt, bottom=2pt]{
    \textbf{\textcolor{azulprofundo}{📚 CLASES DE DOCUMENTO}}
}
\vspace{0.3cm}

\begin{tabular}{@{}ll@{}}
    \toprule
    \textbf{Clase} & \textbf{Uso} \\
    \midrule
    \comando{article} & Artículos, informes cortos \\
    \comando{report} & Informes largos, tesis \\
    \comando{book} & Libros \\
    \comando{beamer} & Presentaciones \\
    \comando{IEEEtran} & Artículos IEEE \\
    \bottomrule
\end{tabular}

\vspace{0.5cm}

% ============================================
% SECCIONES
% ============================================
\tcbox[colback=grisclaro, colframe=azulprofundo, arc=3pt, boxsep=2pt, left=2pt, right=2pt, top=2pt, bottom=2pt]{
    \textbf{\textcolor{azulprofundo}{📑 SECCIONES}}
}
\vspace{0.3cm}

\begin{tabular}{@{}ll@{}}
    \toprule
    \textbf{Comando} & \textbf{Nivel} \\
    \midrule
    \comando{\textbackslash section\{...\}} & Sección \\
    \comando{\textbackslash subsection\{...\}} & Subsección \\
    \comando{\textbackslash subsubsection\{...\}} & Subsubsección \\
    \comando{\textbackslash paragraph\{...\}} & Párrafo \\
    \comando{\textbackslash part\{...\}} & Parte \\
    \bottomrule
\end{tabular}

\vspace{0.5cm}

% ============================================
% FORMATO DE TEXTO
% ============================================
\tcbox[colback=grisclaro, colframe=azulprofundo, arc=3pt, boxsep=2pt, left=2pt, right=2pt, top=2pt, bottom=2pt]{
    \textbf{\textcolor{azulprofundo}{✏️ FORMATO DE TEXTO}}
}
\vspace{0.3cm}

\begin{tabular}{@{}ll@{}}
    \toprule
    \textbf{Comando} & \textbf{Efecto} \\
    \midrule
    \comando{\textbackslash textbf\{...\}} & \textbf{Negrita} \\
    \comando{\textbackslash textit\{...\}} & \textit{Cursiva} \\
    \comando{\textbackslash underline\{...\}} & \underline{Subrayado} \\
    \comando{\textbackslash texttt\{...\}} & \texttt{Monoespaciado} \\
    \comando{\textbackslash textsc\{...\}} & \textsc{Versalitas} \\
    \comando{\textbackslash small} & \small{Pequeño} \\
    \comando{\textbackslash large} & \large{Grande} \\
    \bottomrule
\end{tabular}

\vspace{0.5cm}

% ============================================
% MATEMÁTICAS
% ============================================
\tcbox[colback=grisclaro, colframe=verdeexito, arc=3pt, boxsep=2pt, left=2pt, right=2pt, top=2pt, bottom=2pt]{
    \textbf{\textcolor{verdeexito}{📐 MATEMÁTICAS}}
}
\vspace{0.3cm}

\begin{tabular}{@{}ll@{}}
    \toprule
    \textbf{Comando} & \textbf{Resultado} \\
    \midrule
    \comando{\$...\$} & Modo inline \\
    \comando{\textbackslash[...\textbackslash]} & Modo display \\
    \comando{\textbackslash begin\{equation\}} & Ec. numerada \\
    \comando{\textbackslash frac\{a\}\{b\}} & $\frac{a}{b}$ \\
    \comando{\textbackslash sqrt\{x\}} & $\sqrt{x}$ \\
    \comando{\textbackslash sqrt[n]\{x\}} & $\sqrt[n]{x}$ \\
    \comando{x\_\{i\}} & $x_i$ \\
    \comando{x\^\{n\}} & $x^n$ \\
    \bottomrule
\end{tabular}

\vspace{0.5cm}

% ============================================
% SÍMBOLOS MATEMÁTICOS
% ============================================
\tcbox[colback=grisclaro, colframe=verdeexito, arc=3pt, boxsep=2pt, left=2pt, right=2pt, top=2pt, bottom=2pt]{
    \textbf{\textcolor{verdeexito}{🔣 SÍMBOLOS MATEMÁTICOS}}
}
\vspace{0.3cm}

\begin{tabular}{@{}ll@{}}
    \toprule
    \textbf{Comando} & \textbf{Símbolo} \\
    \midrule
    \comando{\textbackslash alpha} & $\alpha$ \\
    \comando{\textbackslash beta} & $\beta$ \\
    \comando{\textbackslash gamma} & $\gamma$ \\
    \comando{\textbackslash pi} & $\pi$ \\
    \comando{\textbackslash Omega} & $\Omega$ \\
    \comando{\textbackslash infty} & $\infty$ \\
    \comando{\textbackslash sum} & $\sum$ \\
    \comando{\textbackslash int} & $\int$ \\
    \comando{\textbackslash prod} & $\prod$ \\
    \comando{\textbackslash partial} & $\partial$ \\
    \bottomrule
\end{tabular}

\vspace{0.5cm}

% ============================================
% OPERADORES
% ============================================
\tcbox[colback=grisclaro, colframe=verdeexito, arc=3pt, boxsep=2pt, left=2pt, right=2pt, top=2pt, bottom=2pt]{
    \textbf{\textcolor{verdeexito}{⚡ OPERADORES}}
}
\vspace{0.3cm}

\begin{tabular}{@{}ll@{}}
    \toprule
    \textbf{Comando} & \textbf{Símbolo} \\
    \midrule
    \comando{\textbackslash cdot} & $\cdot$ \\
    \comando{\textbackslash times} & $\times$ \\
    \comando{\textbackslash div} & $\div$ \\
    \comando{\textbackslash pm} & $\pm$ \\
    \comando{\textbackslash mp} & $\mp$ \\
    \comando{\textbackslash leq} & $\leq$ \\
    \comando{\textbackslash geq} & $\geq$ \\
    \comando{\textbackslash neq} & $\neq$ \\
    \comando{\textbackslash approx} & $\approx$ \\
    \bottomrule
\end{tabular}

\vspace{0.5cm}

% ============================================
% LISTAS
% ============================================
\tcbox[colback=grisclaro, colframe=azulprofundo, arc=3pt, boxsep=2pt, left=2pt, right=2pt, top=2pt, bottom=2pt]{
    \textbf{\textcolor{azulprofundo}{📋 LISTAS}}
}
\vspace{0.3cm}

\begin{tabular}{@{}ll@{}}
    \toprule
    \textbf{Entorno} & \textbf{Tipo} \\
    \midrule
    \comando{itemize} & Viñetas \\
    \comando{enumerate} & Numerada \\
    \comando{description} & Descriptiva \\
    \bottomrule
\end{tabular}

\vspace{0.3cm}

\textbf{Ejemplo itemize:}
\begin{lstlisting}[language=TeX, basicstyle=\ttfamily\scriptsize]
\begin{itemize}
    \item Elemento 1
    \item Elemento 2
\end{itemize}
\end{lstlisting}

\vspace{0.5cm}

% ============================================
% TABLAS
% ============================================
\tcbox[colback=grisclaro, colframe=azulprofundo, arc=3pt, boxsep=2pt, left=2pt, right=2pt, top=2pt, bottom=2pt]{
    \textbf{\textcolor{azulprofundo}{📊 TABLAS}}
}
\vspace{0.3cm}

\begin{tabular}{@{}ll@{}}
    \toprule
    \textbf{Comando} & \textbf{Función} \\
    \midrule
    \comando{tabular} & Entorno de tabla \\
    \comando{\textbackslash hline} & Línea horizontal \\
    \comando{\&} & Separador de columnas \\
    \comando{\textbackslash\textbackslash} & Salto de fila \\
    \comando{l} & Columna izquierda \\
    \comando{c} & Columna centrada \\
    \comando{r} & Columna derecha \\
    \comando{|} & Línea vertical \\
    \bottomrule
\end{tabular}

\vspace{0.5cm}

% ============================================
% FIGURAS
% ============================================
\tcbox[colback=grisclaro, colframe=azulprofundo, arc=3pt, boxsep=2pt, left=2pt, right=2pt, top=2pt, bottom=2pt]{
    \textbf{\textcolor{azulprofundo}{🖼️ FIGURAS}}
}
\vspace{0.3cm}

\begin{tabular}{@{}ll@{}}
    \toprule
    \textbf{Comando} & \textbf{Función} \\
    \midrule
    \comando{\textbackslash includegraphics[ancho]\{archivo\}} & Insertar imagen \\
    \comando{\textbackslash begin\{figure\}[h]} & Entorno de figura \\
    \comando{\textbackslash caption\{...\}} & Título \\
    \comando{\textbackslash label\{fig:...\}} & Etiqueta \\
    \bottomrule
\end{tabular}

\vspace{0.5cm}

% ============================================
% REFERENCIAS
% ============================================
\tcbox[colback=grisclaro, colframe=azulprofundo, arc=3pt, boxsep=2pt, left=2pt, right=2pt, top=2pt, bottom=2pt]{
    \textbf{\textcolor{azulprofundo}{🔗 REFERENCIAS}}
}
\vspace{0.3cm}

\begin{tabular}{@{}ll@{}}
    \toprule
    \textbf{Comando} & \textbf{Función} \\
    \midrule
    \comando{\textbackslash label\{clave\}} & Crear etiqueta \\
    \comando{\textbackslash ref\{clave\}} & Referenciar \\
    \comando{\textbackslash pageref\{clave\}} & Referencia a página \\
    \comando{\textbackslash cite\{clave\}} & Citar bibliografía \\
    \bottomrule
\end{tabular}

\vspace{0.5cm}

% ============================================
% ECUACIONES DE ELECTRÓNICA
% ============================================
\tcbox[colback=grisclaro, colframe=verdeexito, arc=3pt, boxsep=2pt, left=2pt, right=2pt, top=2pt, bottom=2pt]{
    \textbf{\textcolor{verdeexito}{🔌 ECUACIONES DE ELECTRÓNICA}}
}
\vspace{0.3cm}

\begin{tabular}{@{}ll@{}}
    \toprule
    \textbf{Ecuación} & \textbf{Código} \\
    \midrule
    $V = I \cdot R$ & \comando{V = I \textbackslash cdot R} \\
    $P = V \cdot I$ & \comando{P = V \textbackslash cdot I} \\
    $X_C = \frac{1}{2\pi f C}$ & \comando{X\_C = \textbackslash frac\{1\}\{2\textbackslash pi f C\}} \\
    $X_L = 2\pi f L$ & \comando{X\_L = 2\textbackslash pi f L} \\
    $Z = \sqrt{R^2 + (X_L - X_C)^2}$ & \comando{Z = \textbackslash sqrt\{R\^2 + (X\_L - X\_C)\^2\}} \\
    $f_0 = \frac{1}{2\pi\sqrt{LC}}$ & \comando{f\_0 = \textbackslash frac\{1\}\{2\textbackslash pi\textbackslash sqrt\{LC\}\}} \\
    \bottomrule
\end{tabular}

\vspace{0.5cm}

% ============================================
% COMANDOS PARA MÁRGENES
% ============================================
\tcbox[colback=grisclaro, colframe=azulprofundo, arc=3pt, boxsep=2pt, left=2pt, right=2pt, top=2pt, bottom=2pt]{
    \textbf{\textcolor{azulprofundo}{📏 CONFIGURACIÓN DE PÁGINA}}
}
\vspace{0.3cm}

\begin{tabular}{@{}ll@{}}
    \toprule
    \textbf{Comando} & \textbf{Función} \\
    \midrule
    \comando{\textbackslash usepackage\{geometry\}} & Cargar paquete \\
    \comando{\textbackslash geometry\{margin=1in\}} & Márgenes de 1 pulgada \\
    \comando{\textbackslash geometry\{top=2cm, bottom=2cm\}} & Márgenes personalizados \\
    \comando{\textbackslash pagestyle\{empty\}} & Sin numeración \\
    \comando{\textbackslash pagenumbering\{roman\}} & Numeración romana \\
    \bottomrule
\end{tabular}

\vspace{0.5cm}

% ============================================
% PAQUETES ESENCIALES
% ============================================
\tcbox[colback=grisclaro, colframe=azulprofundo, arc=3pt, boxsep=2pt, left=2pt, right=2pt, top=2pt, bottom=2pt]{
    \textbf{\textcolor{azulprofundo}{📦 PAQUETES ESENCIALES}}
}
\vspace{0.3cm}

\begin{tabular}{@{}ll@{}}
    \toprule
    \textbf{Paquete} & \textbf{Uso} \\
    \midrule
    \comando{amsmath} & Matemáticas avanzadas \\
    \comando{amssymb} & Símbolos matemáticos \\
    \comando{graphicx} & Incluir imágenes \\
    \comando{booktabs} & Tablas profesionales \\
    \comando{geometry} & Márgenes \\
    \comando{hyperref} & Enlaces interactivos \\
    \comando{babel[spanish]} & Idioma español \\
    \bottomrule
\end{tabular}

\end{multicols}

\vspace{0.5cm}

\begin{center}
    \rule{0.8\textwidth}{0.5pt}
    \vspace{0.3cm}
    
    \textbf{Referencia rápida de LaTeX} \\
    Para más información: \texttt{https://github.com/lrocag-dev/UNTELS\_LATEX\_IA}
    \vspace{0.3cm}
\end{center}

\end{document}
